\documentclass[10pt]{article}
\usepackage[utf8]{inputenc}
\usepackage{amsmath}
\usepackage{geometry}
\usepackage{hyperref}
\usepackage{multicol} % Paket untuk membuat banyak kolom
\usepackage{enumitem} % Paket untuk mengatur kerapatan list

% Memperkecil margin dan mengatur jarak antar kolom
\geometry{a4paper, margin=0.6in, columnsep=0.25in} 

% Memangkas spasi kosong berlebih antar baris list
\setlist{nosep, leftmargin=*} 

\title{Instrumentasi dan Pengukuran}
\author{}
\date{}

\begin{document}
\maketitle

\begin{multicols}{2}

\section*{Bab 1: Dasar Instrumentasi}
\label{bab1}

\subsection*{Sistem Akuisisi Data}
Sistem Akuisisi Data secara umum terdiri dari:
\begin{enumerate}
    \item \textbf{Sensor:} Alat pengindera/pengkonversi stimulus fisis.
    \item \textbf{Pengkondisi Sinyal:} Modifikasi sinyal mentah agar bisa diproses lebih lanjut.
    \item \textbf{ADC (Analog to Digital Converter):} Pengubah sinyal fisika (analog) menjadi data biner (digital).
    \item \textbf{Pemroses Sinyal:} Unit mikrokontroler/komputer yang melakukan komputasinya.
    \item \textbf{Aktuator / Penampil (Display):} Penggerak/penyaji \textit{output} ke lingkungan luar.
\end{enumerate}

\subsection*{Terminologi dan Definisi Penting}
\begin{enumerate}
    \item \textbf{Metrologi:} Ilmu pengetahuan yang berhubungan dan mempelajari tentang pengukuran.
    \item \textbf{Instrumentasi:} Bidang ilmu meliputi perancangan, pembuatan, serta penggunaan instrumen untuk keperluan deteksi, pengukuran, serta pengolahan data.
    \item \textbf{Pengukuran:} Rangkaian kegiatan eksperimen guna menentukan nilai kuantitatif sebuah besaran.
    \item \textbf{Ketelitian (Accuracy):} Kemampuan alat ukur untuk memberikan indikasi / kedekatan dengan harga \textbf{nilai sebenarnya}.
    \item \textbf{Ketepatan (Precision):} Kedekatan nilai-nilai pengukuran individual yang didistribusikan sekitar nilai rata-ratanya saat diulang.
    \item \textbf{Sensitivitas (Sensitivity):} Perbandingan antara sinyal keluaran (respon) instrumen terhadap perubahan variabel masukan yang diukur.
    \item \textbf{Repeatabilitas (Repeatability):} Kemampuan alat ukur untuk menunjukkan hasil yang sama dari proses pengukuran yang dilakukan berulang dan identik.
    \item \textbf{Kesalahan (Error):} Penyimpangan variabel yang diukur dari nilai sebenarnya.
    \item \textbf{Resolusi (Resolution):} Perubahan terkecil pada nilai yang diukur dari respon suatu instrumen.
    \item \textbf{Kalibrasi (Calibration):} Serangkaian kegiatan memverifikasi/mencocokkan kebenaran indikasi ukur sebuah alat dengan membandingkannya ke standar ukur yang tertelusuri secara nasional/internasional.
    \item \textbf{Koreksi (Correction):} Suatu harga aljabar yang ditambahkan pada hasil ukur untuk mengkompensasi penambahan kesalahan sistematik.
    \item \textbf{Ketertelusuran (Traceability):} Terkaitnya hasil pengukuran pada standar nasional/internasional melalui peralatan ukur yang kinerjanya diketahui beserta sistem pembuktiannya.
    \item \textbf{Kehandalan (Reliability):} Kesanggupan alat ukur untuk melaksanakan fungsi yang diisyaratkan untuk suatu periode yang ditetapkan.
    \item \textbf{Ketidakpastian Pengukuran (Uncertainty):} Perkiraan atau taksiran rentang dari nilai pengukuran di mana nilai sebenarnya dari besaran obyek yang diukur terletak.
    \item \textbf{Sensor:} Bagian/elemen dari alat ukur untuk mengubah suatu bentuk besaran fisik (sensing element) menjadi suatu \textbf{besaran listrik}.
    \item \textbf{Transduser:} Bagian dari alat ukur yang mengubah suatu bentuk format energi menjadi \textbf{bentuk energi yang lain} sehingga mudah diolah oleh peralatan berikutnya.
    \item \textbf{Rentang ukur (Range):} Besar daerah ukur di antara batas ukur bawah dan batas ukur atas.
    \item \textbf{Jangkauan (Span):} Beda modulus antara dua batas rentang nominal dari alat ukur (Contoh: Bekerja pada rentang -10V sampai 10V, maka Span = 20V).
\end{enumerate}

\noindent\rule{\columnwidth}{0.4pt}

\section*{Bab 2: Alat Ukur Listrik}
\label{bab2}

\subsection*{1. Besaran Listrik Dasar}
\begin{itemize}
    \item \textbf{Arus Listrik (I):} Aliran muatan listrik dalam suatu penghantar, dihitung dengan satuan Ampere (A). Arus mengalir dikarenakan adanya sumber tegangan (voltage) yang terhubung dalam sistem tertutup, yang juga memiliki hambatan (resistansi).
    \item \textbf{Tegangan (E/V):} Besarnya gaya pendorong aliran arus (Electromotive Force), bersatuan Volt (V). Tegangan muncul berkat perbedaan potensial antara dua kutub.
\end{itemize}

\subsection*{2. Galvanometer}
\begin{itemize}
    \item Alat ukur dasar atau komponen pengerak utama di dalam berbagai instrumen ukur listrik tipe jarum penunjuk (analog).
    \item Bekerja berdasarkan prinsip interaksi medan magnet. Arus listrik yang dilewatkan pada sebuah kumparan (coil) di dalam medan magnet permanen akan menghasilkan gaya Lorentz pembentuk momen torsi. 
\end{itemize}

\subsection*{3. Varian Alat Ukur}
\begin{itemize}
    \item \textbf{Amperemeter:} Galvanometer yang dipararelkan dengan tahanan \textit{Shunt} kecil. Harus \textbf{dipasang secara Seri} dengan beban.
    \item \textbf{Voltmeter:} Galvanometer diserikan dengan tahanan \textit{Multiplier} bernilai besar. Harus \textbf{dipasang secara Paralel}.
    \item \textbf{Ohmmeter:} Terdapat tambahan sumber baterai lokal di dalam alat untuk menginjeksi listrik ke komponen yang diuji.
\end{itemize}

\subsection*{4. Multimeter}
Sering disebut AVO-meter (Ampere-Volt-Ohm), menyatukan fungsi Voltmeter, Amperemeter, dan Ohmmeter dalam 1 wadah.

\noindent\rule{\columnwidth}{0.4pt}

\section*{Bab 3: Sensor \& Transduser}
\label{bab3}

\subsection*{Definisi Sensor \& Transduser}
\begin{itemize}
    \item \textbf{Sensor:} Perangkat yang menerima stimulus dan merespon dengan mengeluarkan sinyal elektronik dengan merubah besaran fisis menjadi besaran listrik.
    \item \textbf{Transduser:} Perangkat yang merubah besaran fisis menjadi besaran fisis yang lain.
\end{itemize}

\subsection*{Tipe Konversi \& Sistem}
\begin{enumerate}
    \item \textbf{Sensor Langsung (Direct):} Mengubah stimulus langsung menjadi sinyal listrik lewat 1 efek fisika.
    \item \textbf{Sensor Kompleks:} Melalui lebih dari 1 tahapan konversi transduser.
\end{enumerate}

\subsection*{Klasifikasi Sensor}
\textbf{1. Berdasar Kedudukan Acuan:}
\begin{itemize}
    \item \textbf{Absolut:} Mendeteksi stimulus mandiri pada skala fisis murni independen.
    \item \textbf{Relatif:} Mengukur berdasar kasus khusus / referensi banding gradien tertentu.
\end{itemize}

\textbf{2. Berdasar Ketergantungan Suplai:}
\begin{itemize}
    \item \textbf{Aktif:} Membutuhkan injeksi sumber daya luar (sinyal eksitasi).
    \item \textbf{Pasif:} Secara alami memproduksi energi sinyal listrik tanpa dicatu daya eksternal.
\end{itemize}

\noindent\rule{\columnwidth}{0.4pt}

\section*{Bab 4: Karakteristik Sensor}
\label{bab4}

\subsection*{Karakteristik Statis}
\begin{enumerate}
    \item \textbf{Fungsi Transfer:} Hubungan matematis ideal sinyal input dengan output.
    \begin{itemize}
        \item Linier: $S = a + bs$
        \item Logaritmik: $S = a + b \ln s$
        \item Eksponensial: $S = a e^{bs}$
        \item Pangkat: $S = a_0 + a_1 s^b$
    \end{itemize}
    \textit{(S = Sinyal keluaran, s = Sinyal masukan, b = Sensitivitas)}
    \item \textbf{Span (Range) \& FSO:} Jangkauan dinamis input minimum hingga maksimum. FSO adalah selisih \textit{output} tegangan maksimum dengan minimum.
    \item \textbf{Inakurasi / Error:} 
    \[ \text{Error} = \left(\frac{\text{Deviasi max}}{\text{Span Input}}\right) \times 100\% \]
    \item \textbf{Histeresis:} Ketidakcocokan pembacaan saat grafik gerak naik vs turun.
    \item \textbf{Ketidaklinearan:} Simpangan kurva keluaran terhadap acuan linier idealnya.
    \item \textbf{Saturasi:} Batas tertinggi ukur di mana sensor tidak lagi merespons normal.
    \item \textbf{Pita Mati (Dead Band):} Rentang di mana perubahan masukan belum membangkitkan sinyal \textit{output}.
    \item \textbf{Resolusi:} Perubahan input paling kecil yang sudah dapat dideteksi.
\end{enumerate}

\subsection*{Faktor Dinamis}
\begin{itemize}
    \item \textbf{Eksitasi:} Sinyal listrik pen-trigger sensor.
    \item \textbf{Ketahanan:} Stabilitas \textit{Short Term} (noise) dan \textit{Long Term} (umur bahan).
\end{itemize}

\noindent\rule{\columnwidth}{0.4pt}

\section*{Bab 5: Pengkondisi Sinyal}
\label{bab5}

\subsection*{1. Rangkaian Pembagi Tegangan}
\[ V_{out} = V_{in} \times \left( \frac{R_1}{R_1 + R_2} \right) \]

\subsection*{2. Konverter Arus Ke Tegangan}
\[ V_o = -i \cdot R_F \]

\subsection*{3. Operational Amplifier (Op-Amp)}
\begin{itemize}
    \item \textbf{Inverting:}
    \[ V_{out} = -\left( \frac{R_F}{R_{in}} \right) V_{in} \]
    \item \textbf{Non-Inverting:}
    \[ V_{out} = \left( 1 + \frac{R_F}{R_{in}} \right) V_{in} \]
\end{itemize}

\subsection*{4. Filter Frekuensi (RC Filter)}
Menyaring frekuensi. Titik \textit{Cut-off Frequency} ($f_c$):
\[ f_c = \frac{1}{2 \pi R C} \]

\noindent\rule{\columnwidth}{0.4pt}

\section*{Bab 6: Konversi Sinyal}
\label{bab6}

\subsection*{Format Sinyal}
\begin{itemize}
    \item \textbf{Analog:} Sinyal \textit{kontinu}.
    \item \textbf{Digital:} Sinyal \textit{diskrit}. Dimodulasi dengan \textit{PCM (Pulse Code Modulation)}.
\end{itemize}

\subsection*{Konverter Perantara}
\begin{enumerate}
    \item \textbf{ADC (Analog to Digital):}
    \[ V_{out(Des)} = \left( \frac{V_{in}}{V_{ref}} \right) \times ADC_{\text{max}} \]
    \item \textbf{DAC (Digital to Analog):}
    \[ V_{out} = \left( \frac{\text{In}_{\text{DAC}}}{DAC_{\text{max}}} \right) \times V_{ref} \]
\end{enumerate}

\subsection*{Resolusi ADC Data}
\begin{itemize}
    \item \textbf{Konverter 8-Bit:} $2^8 = 256$ skala (\texttt{0}-\texttt{255}). 
    Resolusi = $V_{ref} / 255$
    \item \textbf{Konverter 10-Bit:} $2^{10} = 1024$ skala (\texttt{0}-\texttt{1023}).
    Resolusi = $V_{ref} / 1023$
\end{itemize}
Semakin besar bit-nya, loncatan nilainya semakin kecil (rapat), sehingga pembacaan lebih presisi.

\end{multicols}

\end{document}